\section{Fazit}
Zentrales Thema dieser Arbeit ist das Exception Handling sowie die dazugehörigen Mechanismen. Ziel ist es, zu zeigen, wie Exception Handling sinnvoll eingesetzt werden kann, um typische Fehler systematisch zu erkennen und zu behandeln. Dazu wurden zunächst die Grundlagen erläutert, insbesondere der Unterschied zwischen Checked und Unchecked Exceptions sowie zwischen Errors.
Durch die konkrete Befassung mit Exception Handling wurden die zentralen Mechanismen sowie deren Bedeutung für eine strukturierte Fehlerbehandlung deutlich. Dabei zeigte sich, dass die Vorgehensweise die Programmiersicherheit und die Wartbarkeit verbessert. In der Praxis trägt sie zu einer höheren Stabilität der Anwendungen bei, da Fehler systematisch identifiziert, verarbeitet und transparent dokumentiert werden.
